\section{The repair theorem}\label{sec:v4-repair}

Throughout this section $(G_1,G_2)$ is a normalized instance
(Definition~\ref{def:v4-normalized}). Every common edge lies on exactly one
common cycle, the common darts at a vertex $v$ are the darts
$c_{\mathrm{out}}(C,v),c_{\mathrm{in}}(C,v)$ of the common cycles $C$ through
$v$, and $U-v$ is connected for every vertex $v$. Every common cycle $C$
carries one fixed traversal, used in both colours; as a cycle of $G_i$ it
lies in a unique block $B_i(C)$ of $G_i$. Let $\mathcal Z_C$ be the set of
components of $U-V(C)$. Blocks are taken with their colour: a block of $G_1$
and a block of $G_2$ are different objects even if they have the same edges.

\begin{definition}\label{def:v4-datum}
A \emph{block datum} consists of a planar embedding $\mathcal E(B)$ of every
block $B$ of $G_1$ and of $G_2$, and a \emph{side assignment} $s$, a bit
$s(C,Z)$ for every common cycle $C$ and $Z\in\mathcal Z_C$. We write
$s(C,x)=s(C,Z)$ for $x\in Z$, and $s(C,Y)$ for the common value on a
connected subgraph $Y$ of $U-V(C)$.
\end{definition}

\begin{theorem}[repair]\label{thm:v4-repair}
A normalized instance $(G_1,G_2)$ has a SEFE if and only if it has a block
datum such that
\begin{enumerate}[label=(R\arabic*),leftmargin=3em]
\item\label{it:v4-R1} at every vertex $v$ there is a cyclic order $\omega_v$
of the common darts at $v$ that restricts, for every block $B$ of either
colour at $v$, to the order of the common darts of $B$ in $\mathcal E(B)(v)$,
and in which, for $i=1,2$, the partition of the common darts by the blocks
of $G_i$ is noncrossing;
\item\label{it:v4-R2} $\side_{\mathcal E(B_i(C))}(C,x)=s(C,x)$ for every
common cycle $C$, colour $i$ and $x\in V(B_i(C))\setminus V(C)$.
\end{enumerate}
Moreover, from such a datum a SEFE is computed in $O(N^2)$ time; its colour
embeddings restrict to $\mathcal E(B)$ on every block $B$, and every vertex
$x\notin V(C)$ lies on side $s(C,x)$ of every common cycle $C$.
\end{theorem}

The orders $\omega_v$ certify compatibility of the block orders; they need
not be the common rotations of the final embedding.

By Lemma~\ref{lem:v4-partitions}, \ref{it:v4-R1} at $v$ is a planarity
condition on a bipartite map on the blocks at $v$, and $\omega_v$ is computed
in linear time.

\emph{Intuition.} The side assignment fixes on which side of every petal at
$v$ every other petal lies, hence the face tree $T_v$ of the common flower at
$v$ (Figure~\ref{fig:v4-flower}). By~\ref{it:v4-R2} each block's own flower
is a subtree of $T_v$; the old order $\omega_v$ only orders the petals around
each node of $T_v$, the contour of this plane tree is the new common
rotation, and the other darts of each block are attached to it.

\begin{lemma}[necessity]\label{lem:v4-necessary}
Let $(\Gamma_1,\Gamma_2)$ be a SEFE, $C$ a common cycle and $Z\in\mathcal
Z_C$. All vertices of $Z$ lie on one side of $C$, the same in both drawings.
The restrictions of $\Gamma_i$ to the blocks and these sides form a block
datum satisfying~\ref{it:v4-R1} and~\ref{it:v4-R2}, with $\omega_v$ the
common rotation at $v$.
\end{lemma}

\begin{proof}
$C$ is the same closed curve in both drawings. The vertices of $Z$ with the
edges of both colours between them form a connected set avoiding this
curve. The common rotation at $v$ restricts to the order of the common
darts of every block, and the partitions are noncrossing by
Lemma~\ref{lem:v4-assembly}(a); \ref{it:v4-R2} holds since sides do not
change under restriction.
\end{proof}

\subsection{Centres and signature trees}\label{sec:v4-sigtree}

A \emph{centre} is a vertex on at least one common cycle; the common cycles
through a centre $v$ are its \emph{petals}, forming the set $F_v$. They
pairwise share only $v$. For a block $B$ of $G_i$ at $v$ let
\[
J_v(B)=\{C\in F_v: B_i(C)=B\};
\]
these are the petals whose darts at $v$ are darts of $B$, so $B$ has a
common dart at $v$ if and only if $J_v(B)\neq\emptyset$. If $C\in F_v\setminus
J_v(B)$, then $V(B)\cap V(C)=\{v\}$, as distinct blocks share at most one
vertex; moreover $C-v$ lies in the component $D_{B_i(C)}\neq D_B$ of $G_i-v$
(Section~\ref{sec:v4-assembly}). So $D_B$ is connected and disjoint from
$V(C)$, and $s(C,B-v)=s(C,D_B)$ is defined.

Fix a side assignment $s$ and a centre $v$. For a petal $E$ and $b\in\Ftwo$
let $q(E,b)\in\Ftwo^{F_v}$ have $q(E,b)_E=b$ and $q(E,b)_C=s(C,E-v)$ for
$C\neq E$ (the path $E-v$ is connected and disjoint from $V(C)$). The
\emph{signature tree} $T_v(s)$ has the distinct vectors $q(E,b)$ as its
vertices and, for every petal $E$, an edge $E$ joining its \emph{ends}
$q(E,0)\neq q(E,1)$. For a connected subgraph $Y$ of $U-v$ disjoint from
all petals let $q(Y)_C=s(C,Y)$.

\begin{lemma}[signature tree]\label{lem:v4-sigtree}
For every side assignment $s$ and centre $v$, $T_v(s)$ is a tree with
$|F_v|$ edges. For every petal $C$ the two components of $T_v(s)-C$ consist
of the vertices with $C$-coordinate $0$ and $1$, and contain $q(C,0)$ and
$q(C,1)$. For every connected subgraph $Y$ of $U-v$ disjoint from all
petals, $q(Y)$ is a vertex of $T_v(s)$.
\end{lemma}

\begin{proof}
\emph{Forest.} The coordinate $C$ changes along the edge $C$ and along no
other edge, so along a cycle of $T_v(s)$ it would change exactly once.
\emph{Connected.} Call petals $E\neq F$ \emph{adjacent} if a path $P$ in
$U-v$ without inner vertices on petals joins $E-v$ to $F-v$. Cutting a path
of the connected graph $U-v$ at its vertices on petals shows that the
adjacency graph is connected. For adjacent $E,F$ and $C\notin\{E,F\}$ the
set $(E-v)\cup P\cup(F-v)$ is connected and avoids $V(C)$, so
$s(C,E-v)=s(C,F-v)$; hence $q(E,s(E,F-v))=q(F,s(F,E-v))$, and the edges
$E,F$ share an end. So $T_v(s)$ is a tree, the coordinate $C$ is constant on
the components of $T_v(s)-C$, and the ends of $C$ lie in different ones.
\emph{Signatures.} A path in $U-v$ from $Y$ to the petals first meets them
at $w\in E-v$; as the path avoids all other petals before $w$,
$q(Y)_C=s(C,E-v)$ for $C\neq E$, and $q(Y)\in\{q(E,0),q(E,1)\}$.
\end{proof}

\begin{figure}[t]
\centering
\begin{tikzpicture}[scale=0.8,
  vtx/.style={circle,fill,inner sep=1.2pt},
  lab/.style={font=\scriptsize}]
\begin{scope}
\node[vtx,label={[lab]below:$v$}] (v) at (0,0) {};
\draw[thick] (0,0) to[out=-40,in=-90] node[vtx,pos=0.55]{} (3,0)
  to[out=90,in=40] node[vtx,pos=0.45]{} (0,0);
\draw[thick] (0,0) to[out=-22,in=-90] node[vtx,pos=0.55]{} (1.4,0)
  to[out=90,in=22] node[vtx,pos=0.45]{} (0,0);
\draw[thick] (0,0) to[out=140,in=90] node[vtx,pos=0.55]{} (-1.9,0)
  to[out=-90,in=-140] node[vtx,pos=0.45]{} (0,0);
\draw[-{Stealth[length=5pt]}] (3,-0.12) -- (3,0.12);
\draw[-{Stealth[length=5pt]}] (1.4,-0.1) -- (1.4,0.12);
\draw[-{Stealth[length=5pt]}] (-1.9,0.12) -- (-1.9,-0.12);
\node[lab] at (3.25,0.35) {$E_1$};
\node[lab] at (1.62,0.3) {$E_2$};
\node[lab] at (-2.15,0.35) {$E_3$};
\node[lab] at (0.7,0) {$001$};
\node[lab] at (2.1,0) {$011$};
\node[lab] at (-1.05,0) {$110$};
\node[lab] at (0.6,1.15) {$111$};
\end{scope}
\begin{scope}[xshift=4.4cm]
\node[vtx,label={[lab]below:$110$}] (a) at (0,0) {};
\node[vtx,label={[lab]below:$111$}] (b) at (1.2,0) {};
\node[vtx,label={[lab]below:$011$}] (c) at (2.4,0) {};
\node[vtx,label={[lab]below:$001$}] (d) at (3.6,0) {};
\draw[thick] (a) -- node[lab,above] {$E_3$} (b)
  -- node[lab,above] {$E_1$} (c) -- node[lab,above] {$E_2$} (d);
\end{scope}
\end{tikzpicture}
\caption{A flower with petals $E_1,E_2,E_3$ at $v$, all traversed
counter-clockwise, so that the inside of a petal is its left side $0$
(left), and its signature tree (right). Vectors list the coordinates of
$E_1,E_2,E_3$; for example $q(E_2,0)=001$ because $E_2-v$ lies inside $E_1$
and outside $E_3$. The counter-clockwise order of the six petal darts at $v$
is the contour of the tree (Lemma~\ref{lem:v4-flower}).}
\label{fig:v4-flower}
\end{figure}

\subsection{Flowers inside blocks}\label{sec:v4-flowers}

A \emph{flower} with centre $v$ is a union of directed cycles, its petals,
pairwise sharing exactly $v$. In a planar drawing of a flower $F$ with petal
set $J$, the \emph{face tree} $S(F)$ has the faces as nodes and, for every
petal $P$, an edge $P$ joining the face on side $0$ of $P$ along the edges of
$P$, the end $(P,0)$, to the face on side $1$, the end $(P,1)$.

\begin{lemma}[flowers]\label{lem:v4-flower}
For a planarly drawn flower $F$ with centre $v$, $S(F)$ is well defined and
is a tree; distinct faces have distinct side vectors. The rotation at $v$,
with $c_{\mathrm{out}}(P,v)$ read as the traversal of $P$ from $(P,1)$ to
$(P,0)$ and $c_{\mathrm{in}}(P,v)$ as the reverse, is the contour of $S(F)$
for exactly one rotation system, and the rotation at a node $t$ is then the
restriction of the rotation at $v$ to the darts $c_{\mathrm{out}}(P,v)$ of
the edges $P$ at $t$.
\end{lemma}

\begin{proof}
By Lemma~\ref{lem:v4-faces} $F$ has $|J|+1$ faces with distinct side
vectors. Points near an edge of $P$ on side $b$ of $P$ lie, for $P'\neq P$,
on the side of $P'$ containing $P-v$; so all edges of $P$ have the same face
on side $b$, and the two ends differ. Every face is bounded by an edge of
some petal, so it is an end and touches $v$. A small counter-clockwise
circle around $v$ crosses $c_{\mathrm{out}}(P,v)$ from side $1$ to side $0$
of $P$ and $c_{\mathrm{in}}(P,v)$ the other way; it is a closed walk in
$S(F)$ that visits every node and uses every edge once in each direction, so
$S(F)$ is connected with $|J|+1$ nodes and $|J|$ edges, a tree. Let
$\rho(t)$ map each edge along which the walk enters $t$ to the edge along
which it leaves next. If $\rho(t)$ had several cycles, splitting $t$
accordingly would disconnect the tree, while the walk is a single orbit of
the successor map of the split graph using all edges. So $\rho$ is a
rotation system with the walk as its contour, and it is determined by the
contour. The last statement is Lemma~\ref{lem:v4-median}(b).
\end{proof}

For a block $B$ of $G_i$ at a centre $v$ with $J_v(B)\neq\emptyset$, the
petals in $J_v(B)$ form a flower $F_B\subseteq B$ drawn by $\mathcal E(B)$;
let $S_B=S(F_B)$ with the rotation system of Lemma~\ref{lem:v4-flower}. For
a face $f$ of $F_B$ let $\varphi_B(f)_C$ be the side of $C$ containing $f$ if
$C\in J_v(B)$, and $s(C,B-v)$ otherwise.

\begin{lemma}[block subtrees]\label{lem:v4-subtree}
Assume~\ref{it:v4-R2}. Then $\varphi_B$ is an isomorphism from $S_B$ onto the
subtree $T_B$ of $T_v(s)$ formed by the edges $J_v(B)$ and their ends, mapping
the end $(E,b)$ to $q(E,b)$; every vertex of $T_B$ has $C$-coordinate
$s(C,B-v)$ for $C\notin J_v(B)$.
\end{lemma}

\begin{proof}
Let $f$ be the end $(E,b)$; then $\varphi_B(f)_E=b$. For $C\in J_v(B)$,
$C\neq E$, points near the edges of $E$ lie on the side of $C$ containing
$E-v$, which is $s(C,E-v)$ by~\ref{it:v4-R2}, since $V(E-v)\subseteq
V(B)\setminus V(C)$ and $B=B_i(C)$. For $C\notin J_v(B)$,
$s(C,B-v)=s(C,E-v)$ as $E-v\subseteq B-v$. So $\varphi_B(f)=q(E,b)$. Every
face is an end, so $\varphi_B$ maps nodes to vertices and the edge $E$ to the
edge $E$; it is injective since distinct faces of $F_B$ have distinct side
vectors over $J_v(B)$. An injective map of a tree that sends edges to edges
is an isomorphism onto its image.
\end{proof}

\subsection{Orders at a centre}\label{sec:v4-centre}

\begin{lemma}[private runs]\label{lem:v4-runs}
Let $\mathcal B$ be a set of pairwise disjoint finite sets with cyclic orders
$\theta_B$, and $\emptyset\ne M\subseteq\bigcup\mathcal B$ a set of
\emph{marked} elements. Let $\omega$ be a cyclic
order of $M$ that restricts to $\theta_B|_{M\cap B}$ for every $B$ and in
which the sets $M\cap B$ are noncrossing. Replace every marked $d\in B$ by
$d$ followed by the \emph{run} of unmarked elements that follow $d$ in
$\theta_B$ before the next marked one, and insert every $B$ with $M\cap
B=\emptyset$, cut anywhere, into the corner after the last element of the
run of some marked $d$ (several sets one after another). The result $\omega^+$ restricts to every $\theta_B$, the
sets $B$ are noncrossing in it, and every unmarked element of a set with a
marked element, or inserted after the run of $d$, lies between $d$ and its
successor in $\omega$.
\end{lemma}

\begin{proof}
The runs of $B$ appear in the order of their marked elements, which is
$\theta_B|_{M\cap B}$, so $\omega^+$ restricts to $\theta_B$; inserted sets
are intervals. Suppose sets $B,B'$ with marked elements interleave via
$a,b,a',b'$. Delete the inserted sets; the runs are intervals, and $a,a'$
lie in different runs since an interval containing both would contain $b$
or $b'$. By Lemma~\ref{lem:v4-intervals}(a) the marked heads of the runs of
$a,b,a',b'$ interleave in $\omega$, a contradiction.
\end{proof}

\begin{lemma}[orders at a centre]\label{lem:v4-centre}
Let a block datum satisfy~\ref{it:v4-R1} and~\ref{it:v4-R2}, and let $v$ be a
centre. There are cyclic orders $\omega^1_v,\omega^2_v$ of the darts of $G_1$
and of $G_2$ at $v$ such that
\begin{enumerate}[label=(\roman*),nosep]
\item $\omega^i_v$ restricts to $\mathcal E(B)(v)$ for every block $B$ of
$G_i$ at $v$, and the dart sets of distinct such blocks do not interleave;
\item $\omega^1_v$ and $\omega^2_v$ restrict to the same order of the common
darts;
\item for every block $B$ of $G_i$ at $v$ and $C\in F_v\setminus J_v(B)$,
$\gap_{\omega^i_v}(C,v,d)=s(C,B-v)$ for every dart $d$ of $B$ at $v$.
\end{enumerate}
\end{lemma}

\begin{proof}
Let $T=T_v(s)$ and $\omega_v$ as in~\ref{it:v4-R1}. For a node $t$ of $T$
with edge set $J_t$ let $\rho_t$ be the restriction of $\omega_v$ to the
darts $c_{\mathrm{out}}(E,v)$, $E\in J_t$, read as an order of $J_t$. Let
$\omega'_v$ be the contour of the plane tree $(T,\rho)$, read as an order of
the common darts as in Lemma~\ref{lem:v4-flower}. Let $B$, $B'$ be distinct
blocks of $G_i$ at $v$ with common darts.

(1) \emph{$\omega'_v$ restricts to $\mathcal E(B)(v)$ on the common darts of
$B$.} By Lemma~\ref{lem:v4-flower} and~\ref{it:v4-R1}, the rotation of
$T_B$ at $t$ transported from $S_B$ by Lemma~\ref{lem:v4-subtree} is the
restriction of $\omega_v$, hence of $\rho_t$. So $T_B$ with the restricted
rotations is $S_B$, and by Lemma~\ref{lem:v4-median}(c) the restriction of
$\omega'_v$ is the contour of $S_B$, which is the order of the petal darts
in $\mathcal E(B)(v)$ (Lemma~\ref{lem:v4-flower}).

(2) \emph{The common darts of $B$ and $B'$ do not interleave in $\omega'_v$.}
The subtrees $T_B,T_{B'}$ are edge-disjoint, so they share at most one node.
If they share none, let $e$ be the edge by which the path from $T_B$ to
$T_{B'}$ leaves $T_B$; by Lemma~\ref{lem:v4-median}(a) the traversals behind
$e$ form an interval containing the darts of $B'$ and none of $B$. If they
share $t$, every edge of $T_B$ is at $t$ or behind an edge of $T_B$ at $t$,
and no edge of $T_{B'}$ is behind one; so the darts of $B$ and $B'$ lie in
the excursions (Lemma~\ref{lem:v4-median}(b)) along $J_v(B)\cap J_t$ and
$J_v(B')\cap J_t$, which do not interleave in $\rho_t$ by~\ref{it:v4-R1}, and
Lemma~\ref{lem:v4-intervals}(a) applies.

(3) \emph{Let $C\in F_v$ and let $T^C_b$ be the component of $T-C$ with
$C$-coordinate $b$. For a common dart $q$ whose traversal ends in $T^C_b$,
$\lambda_{\omega'_v}(C,v,q)=b$.} By Lemma~\ref{lem:v4-median}(a) for the
traversal $c_{\mathrm{out}}(C,v)$, from $q(C,1)$ to $q(C,0)$, the open
interval $(c_{\mathrm{out}}(C,v),c_{\mathrm{in}}(C,v))_{\omega'_v}$ consists
of the traversals of the edges of $T^C_0$; a traversal of an edge of $T^C_b$
ends in $T^C_b$, $c_{\mathrm{out}}(C,v)$ ends in $T^C_0$ and
$c_{\mathrm{in}}(C,v)$ in $T^C_1$.

\emph{Completion.} Fix $i$. A block $Y$ of $G_i$ at $v$ without common darts
has $Y-v$ disjoint from all petals, so $q_Y=q(Y-v)$ is a vertex of $T$
(Lemma~\ref{lem:v4-sigtree}); choose a common dart $d_Y$ whose traversal
ends at $q_Y$. Apply Lemma~\ref{lem:v4-runs} to $\omega'_v$, the rotations
$\mathcal E(B)(v)$ of the blocks of $G_i$ at $v$ and the common darts as
marked elements (there are some, as $v$ is a centre, and each is a dart of a
block of $G_i$ at $v$), inserting every such $Y$ after the run of $d_Y$; let
$\omega^i_v$ be the result. Its hypotheses hold by (1) and (2), which gives
(i), and (ii) holds by construction. For (iii) let $b=s(C,B-v)$. If $B$ has
common darts, $T_B$ avoids the edge $C$ and its vertices have
$C$-coordinate $b$ (Lemma~\ref{lem:v4-subtree}), so $T_B\subseteq T^C_b$; a
common dart $d$ of $B$ is not a dart of $C$, so
$\gap_{\omega^i_v}(C,v,d)=\lambda_{\omega'_v}(C,v,d)=b$ by (3), and every
private dart of $B$ lies in the corner after a common dart of $B$
(Lemma~\ref{lem:v4-runs}). If $B=Y$ has none, its darts lie in the corner
after $d_Y$, whose traversal ends at $q_Y\in T^C_b$, and (3) applies.
\end{proof}

\subsection{Proof of the repair theorem}\label{sec:v4-global}

\begin{proof}[Proof of Theorem~\ref{thm:v4-repair}]
Necessity is Lemma~\ref{lem:v4-necessary}. Let a block datum
satisfy~\ref{it:v4-R1} and~\ref{it:v4-R2}. For $i=1,2$ let $\mathcal E'_i$ be
the rotation system of $G_i$ with rotation $\omega^i_v$
(Lemma~\ref{lem:v4-centre}) at every centre $v$, the rotation of
$\mathcal E(B)$ at every other vertex that lies in one block $B$, and at every
other cutvertex the concatenation of the rotations of its blocks, each cut
anywhere. By Lemma~\ref{lem:v4-centre}(i) and
Lemma~\ref{lem:v4-assembly}(c), $\mathcal E'_i$ is planar and restricts to
$\mathcal E(B)$ on every block $B$ of $G_i$; it is an embedding since $G_i$
is connected. Every common dart is at a centre, so by
Lemma~\ref{lem:v4-centre}(ii) the common rotations agree.

\emph{Sides.} Let $C$ be a common cycle, $x\notin V(C)$ and $B=B_i(C)$. We
show $\side_{\mathcal E'_i}(C,x)=s(C,x)$. If $x\in V(B)$ this
is~\ref{it:v4-R2}, as $\mathcal E'_i$ restricts to $\mathcal E(B)$.
Otherwise the path in the block-cut tree of $G_i$ from $B$ to $x$ leaves $B$
at a cutvertex $w\in V(B)$, and $x$ lies in a component $D$ of $G_i-w$ with
$V(D)\cap V(B)=\emptyset$; $D=D_{B'}$ for a block $B'\neq B$ at $w$.
If $w\notin V(C)$, $D\cup\{w\}$ is connected and avoids $V(C)$, so $x$ has
the side of $w$, which is $s(C,w)=s(C,x)$ by the first case. If $w\in V(C)$,
then $w$ is a centre, $C\in F_w\setminus J_w(B')$, and by
Lemma~\ref{lem:v4-centre}(iii) and Lemma~\ref{lem:v4-sector}
$\side_{\mathcal E'_i}(C,D)=s(C,B'-w)=s(C,x)$, since $D$ is connected in
$U-V(C)$ and contains $B'-w$ and $x$. By Lemma~\ref{lem:v4-charact},
$\mathcal E'_1,\mathcal E'_2$ are the embeddings of a SEFE with the stated
properties.

\emph{Time.} Labelling the components of all $U-V(C)$ takes $O(N^2)$ time.
At a centre $v$ with $k_v=|F_v|\le\delta_v=\deg_{G_1}(v)+\deg_{G_2}(v)$,
$\omega_v$ takes $O(\delta_v)$ time (Lemma~\ref{lem:v4-partitions}),
$T_v(s)$ takes $O(k_v^2)$ by radix sort of the vectors $q(E,b)$, locating the
$q_Y$ takes $O(k_v\delta_v)$ with a binary trie, and the contour and runs
take $O(\delta_v)$; in all $O(N^2)$, and the assembly is linear.
\end{proof}

\subsection{Lobes}\label{sec:v4-lobes}

\begin{lemma}[counting]\label{lem:v4-count}
Let $W$ be a biconnected loopless multigraph with at least three vertices
and $\mathcal C$ a set of cycles of $W$ whose union is a cactus. Then
$\sum_{C\in\mathcal C}\operatorname{cc}(W-V(C))\le2|E(W)|$, where
$\operatorname{cc}$ counts components.
\end{lemma}

\begin{proof}
Each $C\in\mathcal C$ is a block of the cactus $\bigcup\mathcal C$. Root
every component of its block-cut tree at a block, and let $r_C$ be the
parent cutvertex of $C$, or any vertex of $C$ if $C$ is a root. The sets
$V(C)\setminus\{r_C\}$ are pairwise disjoint: a vertex $u$ on several members
is a cutvertex of the cactus, and all blocks at $u$ except its parent are
children of $u$. Let $Y$ be a component of $W-V(C)$. If $Y$ had at most one
neighbour on $C$, this neighbour would be a cutvertex of $W$, or $W$ would be
disconnected; so we may charge $Y$ to an edge $yz$ with $y\in Y$ and $z\in
V(C)\setminus\{r_C\}$. Through its end $z$ an edge is charged at most once,
since $z$ determines $C$ and $y$ determines $Y$. So every edge is charged at
most twice.
\end{proof}

Hence $\sum_C|\mathcal Z_C|\le2|E(U)|$ ($W=U$), and, applied to the
blocks of $G_1,G_2$ containing common cycles, the graphs $B_i(C)-V(C)$ have
at most $2(|E_1|+|E_2|)$ components in all.
